

% ==================================================
% Judul halaman
% ==================================================
\chapter*{\MakeUppercase{Lembar Pengesahan}}

% Menambahkan Lembar Pengesahan ke Daftar Isi
\addcontentsline{toc}{chapter}{\MakeUppercase{Lembar Pengesahan}}

\begin{center}

% Jenis dokumen
\large\MakeUppercase{\textbf{Tugas Akhir}}

\vspace{2em}

% ==================================================
% Judul Tugas Akhir dalam Bahasa Indonesia
% ==================================================
\large\MakeUppercase{\textbf{Judul dalam Bahasa Indonesia}}\\
% Contoh:
% Sistem Deteksi Burung Berbasis TinyML untuk Monitoring Keanekaragaman Hayati

\vspace{2em}

% ==================================================
% Judul Tugas Akhir dalam Bahasa Inggris
% ==================================================
\large\MakeUppercase{\textbf{\textit{Title in English}}}\\
% Example:
% TinyML-Based Bird Detection System for Biodiversity Monitoring

\vspace{2em}

% Keterangan pengesahan
\normalsize \textbf{Telah disetujui dan disahkan sebagai Buku Tugas Akhir}\\

% ==================================================
% Nama program studi
% ==================================================
\textbf{Program Studi Teknik Fisika}\\
% Contoh:
% Program Studi Teknik Elektro
% Program Studi Teknik Telekomunikasi

% Nama fakultas
\textbf{Fakultas Teknik Elektro}\\

% Nama universitas
\textbf{Universitas Telkom}\\

\vspace{2em}

\textbf{Disusun oleh:}\\

% ==================================================
% Nama lengkap mahasiswa
% ==================================================
\MakeUppercase{\textbf{NAMA PENULIS}}\\
% Contoh:
% AHMAD QURTHOBI

% ==================================================
% Nomor Induk Mahasiswa (NIM)
% ==================================================
\MakeUppercase{\textbf{NIM PENULIS}}\\
% Contoh:
% 120220001

\vfill

% ==================================================
% Tanggal pengesahan
% Dapat menggunakan \today atau ditulis manual
% ==================================================
\textbf{Bandung, \today}\\
% Contoh manual:
% \textbf{Bandung, 25 Juni 2026}\\

\vspace{1em}

\begin{tabular}{ccc}

    % ==================================================
    % Nama pembimbing
    % ==================================================
    \textbf{Pembimbing I} && \textbf{Pembimbing II} \\

    % Ruang untuk tanda tangan
     && \\
     && \\
     && \\
     && \\
     && \\
     && \\
     && \\

    % ==================================================
    % Nama lengkap dan gelar pembimbing
    % ==================================================
    \MakeUppercase{\textbf{Nama Pembimbing 1}}
    &
    \hfill
    &
    \MakeUppercase{\textbf{Nama Pembimbing 2}}
    \\
    
    % Contoh:
    % Dr. Ir. Budi Santoso, S.T., M.T.
    % Dr. Ani Wijayanti, S.Si., M.T.

    % ==================================================
    % NIP pembimbing
    % ==================================================
    \MakeUppercase{\textbf{NIP. }}
    &
    \hfill
    &
    \MakeUppercase{\textbf{NIP. }}
    \\

    % Contoh:
    % NIP. 197801012005011001

\end{tabular}

\end{center}
